\chapter{Metrik Tamlık ve Kompaktlık}

Metrik tamlık (completeness), bir metrik uzayda Cauchy dizilerinin yakınsayıp
yakınsamadığını araştıran temel bir kavramdır.

%-------------------------------------------------------------------
\section{Konu}

\subsection{Cauchy Dizisi ve Tamlık}

\begin{definition}[Cauchy Dizisi]
$(M, d)$ metrik uzayında $\{x_n\}$ dizisi \textbf{Cauchy} ise:
\[
\forall\varepsilon>0,\;\exists N \in \mathbb{N}:\;
m,n \geq N \Rightarrow d(x_m, x_n) < \varepsilon.
\]
\end{definition}

\begin{definition}[Tam Metrik Uzay]
$(M, d)$ \textbf{tam} ise: Her Cauchy dizisi $M$'de bir limite yakınsır.
\end{definition}

\begin{figure}[ht]
\centering
\begin{tikzpicture}[font=\small]
  % eksen
  \draw[-{Stealth[length=2.4mm]}] (-0.4,0) -- (9.2,0) node[right] {$n$};
  % limit cizgisi
  \draw[gray!55, dashed] (0,2.1) -- (8.6,2.1) node[right, black] {$L$};
  % Cauchy dizisinin terimleri: L'ye yaklasan, salinim azalan
  \foreach \n/\y in {0/3.5, 1/0.8, 2/3.0, 3/1.2, 4/2.6, 5/1.7, 6/2.35, 7/1.9, 8/2.18}
    \fill[blue!70!black] (\n,\y) circle (1.6pt);
  % kuyrugun capi -> 0 vurgusu (saga dogru daralan bant)
  \fill[orange!18] (5.5,1.55) -- (8.5,2.0) -- (8.5,2.2) -- (5.5,1.85) -- cycle;
  \node[orange!75!black, font=\scriptsize] at (7.0,1.15) {kuyruğun çapı $\to 0$};
  % ilk terimler arasi buyuk mesafe
  \draw[<->, red!70!black] (0,3.5) -- node[right, font=\scriptsize] {büyük} (0,0.8);
  \node[font=\scriptsize] at (1.6,-0.5) {$x_n$ terimleri};
  \node[blue!70!black] at (7.3,3.4) {$\{x_n\}$ Cauchy};
\end{tikzpicture}

\caption{Cauchy dizisinin terimleri birbirine yaklaşır: kuyruğun çapı sıfıra gider.}
\end{figure}

\begin{sezgi}
Bir Cauchy dizisini, ilerledikçe birbirine sokulan bir terim kalabalığı gibi
düşünün. İlk terimler birbirinden uzaktayken, yeterince ileri gidildiğinde tüm
terimler bir noktada toplanır: ``kuyruğun çapı''
$\bigl(\sup_{m,n\geq N} d(x_m, x_n)\bigr)$ sıfıra iner. \textbf{Tam} uzayda bu
kalabalığın sıkıştığı yerde gerçekten bir nokta (limit) durur; tam olmayan
uzayda ise o yerde bir \emph{boşluk} olabilir.
\end{sezgi}

\subsection{Tamamen Sınırlılık}

$A \subseteq M$ \textbf{tamamen sınırlı} ise: her $\varepsilon>0$ için $A$'nın
sonlu bir $\varepsilon$-net kümesi vardır.
\[
S = \{s_1,\ldots,s_n\} \subseteq M,\quad
\forall x \in A\;\exists i:\; d(x, s_i) < \varepsilon.
\]

\subsection{Metrik Kompaktlık Karakterizasyonu}

\[
(M, d) \text{ tam} \wedge \text{tamamen sınırlı}
\Longleftrightarrow
(M, d) \text{ kompakt.}
\]

\begin{table}[h]
\centering
\begin{tabular}{llll}
\toprule
\textbf{Uzay} & \textbf{Tam?} & \textbf{Tam. sınırlı?} & \textbf{Kompakt?} \\
\midrule
$\mathbb{R}$ & \checkmark & $\times$ & $\times$ \\
$[0,1]$ & \checkmark & \checkmark & \checkmark \\
$(0,1)$ & $\times$ & \checkmark & $\times$ \\
$\mathbb{Q}$ & $\times$ & $\times$ & $\times$ \\
Sonlu metrik & \checkmark & \checkmark & \checkmark \\
\bottomrule
\end{tabular}
\caption{Tamlık ve kompaktlık örnekleri}
\end{table}

\begin{figure}[ht]
\centering
\begin{tikzpicture}[font=\small]
  % --- Q dogrusu (ustte): limit yok ---
  \draw[-{Stealth[length=2.4mm]}] (-0.3,2.4) -- (6.6,2.4) node[right] {$\mathbb{Q}$};
  \draw[red!70!black, dashed] (4.2,2.05) -- (4.2,2.75);
  \node[red!70!black, font=\scriptsize] at (4.2,3.05) {$\sqrt{2}\notin\mathbb{Q}$};
  \foreach \x in {1.0, 2.0, 2.8, 3.4, 3.8, 4.0}
    \fill[blue!70!black] (\x,2.4) circle (1.5pt);
  \node[blue!70!black, font=\scriptsize] at (1.7,2.0) {$x_n\to\sqrt{2}$};
  \node[red!70!black] at (5.7,2.95) {limit \emph{yok}};

  % --- R dogrusu (altta): limit var ---
  \draw[-{Stealth[length=2.4mm]}] (-0.3,0.4) -- (6.6,0.4) node[right] {$\mathbb{R}$};
  \fill[green!45!black] (4.2,0.4) circle (2.2pt);
  \node[green!45!black, font=\scriptsize] at (4.2,0.05) {$\sqrt{2}$};
  \foreach \x in {1.0, 2.0, 2.8, 3.4, 3.8, 4.0}
    \fill[blue!70!black] (\x,0.4) circle (1.5pt);
  \node[green!45!black] at (5.7,0.95) {limit \emph{var}};

  % ayrim
  \node[font=\scriptsize] at (3.1,1.4) {aynı Cauchy dizisi};
\end{tikzpicture}

\caption{$\mathbb{Q}$'da $\sqrt{2}$'ye yakınsayan Cauchy dizisinin limiti yoktur;
aynı dizi $\mathbb{R}$'de yakınsar.}
\end{figure}

\begin{karsiornek}
Rasyonel sayılar $\mathbb{Q}$, Öklid metriğiyle \textbf{tam değildir}.
$x_1 = 1,\ x_2 = 1.4,\ x_3 = 1.41,\ x_4 = 1.414,\ \ldots$ dizisi $\sqrt{2}$'nin
ondalık açılımıdır: her terim rasyoneldir ve dizi Cauchy'dir
($d(x_m, x_n) \to 0$). Ama limit $\sqrt{2} \notin \mathbb{Q}$ olduğundan dizi
$\mathbb{Q}$ içinde \emph{hiçbir} noktaya yakınsamaz. Aynı dizi $\mathbb{R}$'de
$\sqrt{2}$'ye yakınsar --- fark uzayda ``boşluk'' olup olmamasıdır.
\end{karsiornek}

\paragraph{Tamlama (completion).}
Her metrik uzay $(M, d)$ bir \textbf{tam} uzayın $(\widehat{M}, \widehat{d})$
yoğun alt-uzayı olarak gömülebilir; $\widehat{M}$'ye $M$'nin \emph{tamlaması}
denir. Sezgisel olarak tamlama, eksik limitleri (``boşlukları'') ekleyerek uzayı
kapatır: $\mathbb{R} = \widehat{\mathbb{Q}}$ bu inşanın klasik örneğidir.

\begin{figure}[ht]
\centering
\begin{tikzpicture}[font=\small,
    kutu/.style={draw, rounded corners, align=center, inner sep=8pt}]
  % delikli uzay
  \node[kutu, fill=orange!8] (q) at (0,0)
    {$\mathbb{Q}$ (tam değil)\\[3pt] \emph{boşluklar} var};
  % delikleri temsil eden noktalar + bosluklar
  \foreach \x in {-0.9, -0.3, 0.4} \fill[blue!70!black] (\x,-0.95) circle (1.3pt);
  \draw[red!70!black, thick] (-0.6,-0.95) ++(-0.05,0.18) -- ++(0.1,-0.36);
  \draw[red!70!black, thick] (0.05,-0.95) ++(-0.05,0.18) -- ++(0.1,-0.36);
  \node[red!70!black, font=\scriptsize] at (0,-1.45) {eksik limitler};

  % tam uzay
  \node[kutu, fill=green!8] (r) at (7.6,0)
    {$\mathbb{R}=\widehat{\mathbb{Q}}$ (tam)\\[3pt] her Cauchy yakınsar};

  % tamlama oku (iki kutu arasi)
  \draw[-{Stealth[length=3mm]}, thick]
    (q.east) -- node[above] {tamlama $\widehat{\;\cdot\;}$}
                node[below, font=\scriptsize] {boşlukları doldur} (r.west);

  \foreach \x in {6.7, 7.3, 7.6, 7.9, 8.5} \fill[green!45!black] (\x,-0.95) circle (1.3pt);
  \node[green!45!black, font=\scriptsize] at (7.6,-1.45) {süreklilik};
\end{tikzpicture}

\caption{Tamlama: $\mathbb{Q}$'nun boşlukları doldurularak tam uzay $\mathbb{R}$
elde edilir.}
\end{figure}

%-------------------------------------------------------------------
\section{Teoremler}

\begin{theorem}[Heine-Borel Metrik Genelleştirmesi]
Metrik uzayda kompaktlık $\Leftrightarrow$ tamlık $\wedge$ tamamen sınırlılık.
\end{theorem}

\begin{theorem}[Banach Sabit-Nokta Teoremi]
$(M, d)$ tam metrik uzay; $T: M \to M$ büzülme ($K < 1$)
$\Rightarrow$ $T$'nin eşsiz sabit noktası $x^*$ vardır: $T(x^*) = x^*$.
\end{theorem}

\begin{proof}[İspat eskizi]
$K < 1$ büzülme sabiti olsun. Herhangi bir $x_0$ alıp $x_{n+1} = T(x_n)$
iterasyonunu kurun. Büzülmeden
$d(x_{n+1}, x_n) \leq K\, d(x_n, x_{n-1}) \leq K^n\, d(x_1, x_0)$. Üçgen
eşitsizliği ve geometrik seri ile
$d(x_{n+m}, x_n) \leq \tfrac{K^n}{1-K}\, d(x_1, x_0)$; $K < 1$ olduğundan bu
ifade $n \to \infty$ iken sıfıra gider, yani $\{x_n\}$ Cauchy'dir.
\textbf{Tamlık} kullanılarak dizi bir $x^*$'ye yakınsar. $T$ sürekli (Lipschitz)
olduğundan $T(x^*) = T(\lim x_n) = \lim x_{n+1} = x^*$: $x^*$ sabit noktadır.
\emph{Eşsizlik:} $T(p) = p$, $T(q) = q$ olsaydı
$d(p,q) = d(Tp, Tq) \leq K\, d(p,q)$, $K < 1$ ile bu ancak $d(p,q) = 0$, yani
$p = q$ ise mümkündür.
\end{proof}

\begin{theorem}[Tam Uzayın Kapalı Alt-Uzayı Tamdır]
$(M, d)$ tam ve $A \subseteq M$ kapalı ise $(A, d)$ de tamdır.
\end{theorem}

\begin{proof}[İspat eskizi]
$\{a_n\} \subseteq A$ bir Cauchy dizisi olsun. $\{a_n\}$ aynı zamanda $M$ içinde
Cauchy'dir ve $M$ tam olduğundan bir $a \in M$ limitine yakınsar. $A$
\textbf{kapalı} olduğundan kendi limit noktalarını içerir; $a_n \to a$ ve
$a_n \in A$ olması $a \in \overline{A} = A$ demektir. Böylece her Cauchy dizisi
$A$ içinde bir limite yakınsar: $A$ tamdır.
\end{proof}

\begin{theorem}[Baire Kategorisi Teoremi]
Tam metrik uzay 1.\ kategoriden değildir.
\end{theorem}

%-------------------------------------------------------------------
\section{Algoritmalar}

\subsection{Sonlu Metrikte is\_complete ve is\_totally\_bounded}

Sonlu metrik uzayda her Cauchy dizisi eninde sonunda sabit: trivially tam.
Taşıyıcı kendisi $\varepsilon$-net: trivially tamamen sınırlı.
Her ikisi: $O(1)$.

\subsection{metric\_compactness\_check}

\begin{algorithm}
\caption{MetricCompactnessCheck$(M, d)$}
\begin{algorithmic}[1]
\State complete $\leftarrow$ is\_complete$(M, d)$
\State totally\_bounded $\leftarrow$ is\_totally\_bounded$(M, d)$
\State \Return complete $\wedge$ totally\_bounded
\end{algorithmic}
\end{algorithm}

%-------------------------------------------------------------------
\section{pytop API}

\begin{lstlisting}[language=Python]
from pytop.metric_completeness import (
    is_complete,
    is_totally_bounded,
    metric_compactness_check,
    analyze_metric_completeness,
)
\end{lstlisting}

%-------------------------------------------------------------------
\section{Örnekler}

\subsection*{Örnek 5.1 — Sonlu Metrik: Tam ve Kompakt}

\begin{lstlisting}[language=Python]
points = ['A', 'B', 'C', 'D']
dist = {(a, b): (0 if a == b else 1)
        for a in points for b in points}
fms = FiniteMetricSpace(carrier=tuple(points), distance=dist)
print("is_complete:", is_complete(fms).status)
print("is_totally_bounded:", is_totally_bounded(fms).status)
mc = metric_compactness_check(fms)
print("metric_compactness:", mc.status, mc.value)
\end{lstlisting}

\begin{verbatim}
is_complete: true
is_totally_bounded: true
metric_compactness: true True
\end{verbatim}

\subsection*{Örnek 5.4 — analyze\_metric\_completeness}

\begin{lstlisting}[language=Python]
r = analyze_metric_completeness(fms)
print("status:", r.status)
print("value:", r.value)
\end{lstlisting}

\begin{verbatim}
status: true
value: {'is_complete': True, 'is_totally_bounded': True, 'metric_compact': True}
\end{verbatim}

\subsection*{Örnek 5.6 — Rasyoneller $\mathbb{Q}$: Tam Olmayan Uzay (Karşı-Örnek)}

\texttt{rationals\_metric()} sembolik bir uzaydır; \texttt{is\_complete} ondalık
taramayla karara bağlanamadığı için \texttt{unknown} döner --- fakat uzayın
\textbf{etiketleri} (tags) $\mathbb{Q}$'nun tam olmadığını kayıt altına alır
(\texttt{'not\_complete'}).

\begin{lstlisting}[language=Python]
from pytop import rationals_metric
from pytop.metric_completeness import is_complete

qm = rationals_metric()
res = is_complete(qm)
print("is_complete:", res.status)
print("'not_complete' tag:", 'not_complete' in qm.tags)
print("justification:", res.justification[0])
\end{lstlisting}

\begin{verbatim}
is_complete: unknown
'not_complete' tag: True
justification: Completeness has exact support only for explicit finite metric spaces.
\end{verbatim}

\subsection*{Örnek 5.7 — Banach ve Sabit-Nokta Profilleri}

\texttt{get\_fixed\_point\_profiles()} sabit-nokta teorisinin küratörlü
profillerini verir; \texttt{fixed\_point\_stability\_summary()} bunları
kararlılığa göre gruplar. Banach büzülme teoremi \textbf{çekici (attracting)}
sabit nokta profiline karşılık gelir.

\begin{lstlisting}[language=Python]
from pytop import get_fixed_point_profiles, fixed_point_stability_summary

profiles = get_fixed_point_profiles()
print("profil sayisi:", len(profiles))
for p in profiles:
    print(f"  {p.key:24s}: {p.stability}")

summary = fixed_point_stability_summary()
print("stable:", summary['stable'])
\end{lstlisting}

\begin{verbatim}
profil sayisi: 5
  attracting_fixed_point  : stable
  repelling_fixed_point   : unstable
  neutral_fixed_point     : neutral
  brouwer_fixed_point     : not_applicable
  periodic_point_n        : not_applicable
stable: ['attracting_fixed_point']
\end{verbatim}

%-------------------------------------------------------------------
\section{Alıştırmalar}

\subsection*{Kodlama}

\begin{enumerate}[label=K\arabic*.]
\item Rasyonel sayılar $\mathbb{Q}$ (sembolik) için \texttt{is\_complete} sonucunu
      kontrol edin.
\item Kendi 4-noktalı metrik uzayınızı oluşturun ve \texttt{metric\_compactness\_check}
      uygulayın.
\item \texttt{analyze\_metric\_completeness} çalıştırın ve \texttt{value} sözlüğünü inceleyin.
\item \texttt{rationals\_metric()} oluşturun; \texttt{is\_complete(...).status} ile
      \texttt{'not\_complete'} etiketinin uzayın \texttt{tags} kümesinde olup olmadığını
      yazdırın. Statü neden \texttt{unknown}, etiket neden \texttt{True}?
\item \texttt{get\_fixed\_point\_profiles()} listesini gezin ve
      \texttt{stability == 'stable'} olan profilin \texttt{key} alanını yazdırın.
      Bu profil hangi teoreme (Banach mı?) karşılık gelir?
\end{enumerate}

\subsection*{Teori}

\begin{enumerate}[label=T\arabic*.]
\item Banach sabit-nokta teoremini sözlü olarak açıklayın ve bir uygulama verin.
\item Kapalı $[0,1]$'in tam, açık $(0,1)$'in tam olmadığını gösterin.
      (Hint: $1/n$ dizisi Cauchy ama $0 \notin (0,1)$.)
\item ``Tam uzayın kapalı alt-uzayı tamdır'' teoremini ispatlayın. Ayrıca tam bir
      uzayda \textbf{açık} bir alt-uzayın tam olması gerekmediğini bir örnekle
      gösterin. (Hint: $\mathbb{R}$ tam, ama $(0,1) \subseteq \mathbb{R}$ açık ve
      tam değil.)
\item $\mathbb{Q}$'nun tam olmadığını $\sqrt{2}$'nin ondalık açılımı üzerinden
      ispatlayın: dizinin Cauchy olduğunu, fakat $\mathbb{Q}$ içinde limitinin
      bulunmadığını gösterin. Tamlama kavramıyla bu ``boşluğun'' nasıl
      kapatıldığını açıklayın.
\end{enumerate}
